
Experiments address three research questions:

\begin{itemize}
\item \textbf{RQ1:} Do the three UQ methods produce measurably different AUROC on HaluEval-QA with LLaMA-2-7B-chat under matched conditions?
\item \textbf{RQ2:} Is the correlation structure between token entropy and semantic entropy consistent with functional NLI filtering?
\item \textbf{RQ3:} Does deberta-large-mnli successfully aggregate stochastic responses on HaluEval-QA short factual QA responses?
\end{itemize}

\begin{table}[htbp]
\centering
\begin{tabular}{ll}
\toprule
Property & Value \\
\midrule
Dataset & HaluEval-QA, 2,000 stratified examples (1,000 hallucinated, 1,000 factual) \\
Model & LLaMA-2-7B-chat (meta-llama/Llama-2-7b-chat-hf), float16 \\
GPU & NVIDIA H100 NVL (single GPU) \\
Stochastic samples N & 5 (temperature = 1.0) \\
Greedy temperature & 0.0 \\
max_new_tokens & 256 \\
NLI model & microsoft/deberta-large-mnli \\
Bootstrap resamples & 1,000 (seed = 42) \\
Bonferroni-corrected α & 0.0167 (k = 3) \\
Min AUROC gap (gate) & 0.05 \\
Random seed & 42 \\
\bottomrule
\end{tabular}
\end{table}

Three sub-hypotheses were executed in sequence, each building on outputs from the prior:

\begin{itemize}
\item \textbf{H-E1 (EXISTENCE, MUST_WORK):} Tests whether at least one UQ method pair exhibits Δ AUROC ≥ 0.05 with non-overlapping 95\% CIs on HaluEval-QA with LLaMA-2-7B-chat.
\item \textbf{H-M1 (MECHANISM, MUST_WORK):} Tests whether Pearson r(TE, SE) < 0.9, using H-E1 inference outputs without new LLM inference.
\item \textbf{H-M2 (MECHANISM, SHOULD_WORK):} Tests whether the NLI aggregation rate — fraction of examples with cluster count < N=5 — is at least 0.50, using H-E1/H-M1 cluster data.
\end{itemize}

A fourth sub-hypothesis H-M3 — testing cross-model ranking stability with Mistral-7B-Instruct (Spearman ρ ≥ 0.8) — was planned but not executed. Results for H-M3 are therefore absent from this paper.

